\section{Classe JsonParser.java}
    \subsection{Ruolo}

    \subsection{Codice originale}
    \inputminted[linenos, frame=lines]{java}{../old/JsonParser.java}

    \subsection{Analisi di Clean Code}
        \subsubsection{Nomi}
        I nomi di variabili risultano significativi, con gli attributi della classe che rispecchiano ciò che rappresentano, con alcune notevoli eccezzioni:
        \begin{itemize}
            \item \textit{root} per quanto abbastanza esplicativo, non rappersenta il root dell'URL ma bensì il path dell'API, \textit{endpoint} o \textit{path} sarebbero pià esplicativi;
            \item \textit{eventsList} riflette il contenuto della variabile, ma non è necessario specificare che si tratti di una lista, \textit{events} sarebbe sufficiente;
            \item \textit{eventObj} riflette il contenuto ma non specifica che si tratti di un oggetto JSON, \textit{eventJson} potrebbe risultare più chiaro.
        \end{itemize}

        Per le funzioni, invece, i nomi risultano poco chiari e spesso non riflettono ciò che la funzione svolge realmente:
        \begin{itemize}
            \item \textit{jsonParse()} non riflette il compito della funzione, che svolge molte altre operazioni oltre al semplice parsing di JSON come la creazione dell'URL, 
            la gestione della connessione e infine il parsing del JSON e gestione degli errori.
        \end{itemize}

        Anche la classe risulta malamente nominata, \textit{JsonParser} suggerisce una classe che esegue il parsing del JSON, in realtà questa si occupa anche di stabilire la connessione e gestirla.
    
        \subsubsection{Responsabilità}
        Come anticipato la funzione \textit{jsonParse()} viola il \textbf{SRP (Single Responsibility Principle)} e il \textbf{CQS (Command Query Separation)}, occupandosi di diversi compiti 
        quando potrebbe facilmente essere divisa in sotto funzioni: 
        \begin{itemize}
            \item creare l'URL della API da chiamare;
            \item scrivere sul Log informazioni di debug;
            \item gestire la connessione HTTP verso la API;
            \item effettuare la chiamata alla API di TicketMaster;
            \item gestire la logica di parsing della risposta JSON;
            \item gestire la logica di eventuali errori di connessione o risposta;
        \end{itemize}

        Questi compiti dovrebbero essere estratti in funzioni separate, migliorando leggibilità e manutenzione del codice, riducendo contemporaneamente la complessità della funzione.

        In generale tutta la classe dovrebbe modificata e divisa in due classi distinte:
        \begin{itemize}
            \item una che si occupi della gestione della connessione alla API di TicketMaster per recuperare i dati (gli eventi in questo caso);
            \item una che si occupi di analizzare e parsare questi dati in formato JSON, estraendo le informazioni richieste.
        \end{itemize}

        \subsubsection{Commenti}
        Dal punto di vista dei commenti, la classe ne presenta un numero eccessivo, alcuni commenti sono utili e chiari, la maggior parte invece interrompe il flusso del codice 
        spiegando concetti ovvi e non fornendo informazioni utili.

        La spiegazione degli attributi della classe è superflua, il nome e il contesto dovrebbero essere sufficienti a comprendere cosa rappresentano.

        In particolare la funzione \textit{jsonParse()} presenta un commento generale che spiega in maniera poco chiara la sua funzionalità, e numerosi commenti all'interno che 
        separano i blocchi di codice per funzionalità differenti, spesso spiegando banalità o fornendo informazioni poco utili.

        SOno presenti anche commenti per codice ormai rimosso, come mostrato alla riga 40.

        \subsubsection{Hardocoding}
        La classe possiede un paio di problemi di embedding, in particolare la \textit{chiave API} è hardcoded in chiaro come attributo statico della classe, 
        ciò rappresenta un problema di sicurezza e manutenzione in caso di futuri cambiamenti.

        Un'altra stringa che risulta essere hardcoded è l'URL base della API, che dovrebbe essere spostata in una costante o variabile globale così come la stringa \textit{\textunderscore embedded}, 
        necessaria per eseguire il parsing della risposta JSON. 

        \subsubsection{Gestione degli errori}
        La gestione degli errori è abbastanza approssimativa; spesso si cattura l'eccezione senza avere comportamenti condizionati al tipo di errore e non si comunica all'utente l'avvenuto problema,
        lasciandolo ad aspettare un evento che non si verificherà mai.
        
    \subsection{Code Refactoring}
